\begin{textsample}{POS Dim 8 – global\_north\_2024\_10 – Score 4.00 – global\_north\_2024\_10/115548283.txt}  \label{ex:f8_pos_249}
In its letter to Education Minister Jill Dunlop, the TDSB asks the province to develop a clear definition of hate that tackles growing antisemitism and other forms of hate and provide resources and training to teachers to help navigate geopolitical conflicts affecting students.
In its letter to Education Minister Jill Dunlop, the TDSB asks the province to develop a clear definition of hate that tackles growing antisemitism and other forms of hate and provide resources and training to teachers to help navigate geopolitical conflicts affecting students.
The Toronto District School Board is asking for help in handling the divisions world conflicts are causing in some classrooms, as several other Ontario boards grapple with similar challenges."Geopolitical conflicts are not \textbf{necessarily} things that are just happening somewhere far away -- they have real consequences on many families in our system,"said Trustee Neethan Shan, acting chair of the TDSB, adding students may have relatives in conflict zones, be impacted by the media \textbf{coverage} or affected by increasing polarization."There ' s pain, there ' s struggle and they ' re trying to heal @ @ @ @ @ @ @ @ @ @ educators to have the necessary training, tools, resources to be able to affirm the different identities and be supportive of students in a way that is not \textbf{exclusive} to one community,"Shan said.
ARTICLE CONTINUES BELOW The board ' s requests for help from the province and Ontario College of Teachers, which regulates the teaching profession, were made as the crisis in the Middle East widens and in the lead-up to the one-year anniversary of the Oct. 7 Hamas attack on Israel and the start of the war in Gaza.
The impact of that conflict has been far-reaching, with incidents of hate on the rise here, impacting both broader communities and schools.
While no recent data is available, the TDSB says it has seen a rise in antisemitism, Islamophobia and anti-Palestinian racism since last fall."We welcome the fact that TDSB has brought this issue to the forefront, as many other boards are working on addressing this growing need as well,"said Kathleen Woodcock, president of the Ontario @ @ @ @ @ @ @ @ @ @ heard from Palestinian parents who say their children feel silenced and are concerned about incidents such as students being sanctioned for wearing kaffiyehs or clothing with"Free Palestine."Student-led walkouts were organized at some Toronto schools in solidarity with people in Gaza.
Meanwhile, Jewish parents and students have also spoken out about racial slurs, kids performing Nazi salutes, antisemitic graffiti and Holocaust denial.
ARTICLE CONTINUES BELOW ARTICLE CONTINUES BELOW Toronto Trustee Shelley Laskin, who introduced the motion calling on the board to reach out for help, said some people are"concerned and hiding their identities so that they don't face potential issues.
And that really concerns me because kids and staff should be comfortable to be their whole selves in schools.
I ' m hearing more of that than actual incidents."The Toronto board also came under fire last month after students from about 15 schools attended a community event to support the Grassy Narrows First Nation and their efforts to address mercury contamination.
Parents were told it @ @ @ @ @ @ @ @ @ @ n't participate in the rally.
But many students and teachers joined the solidarity march for Indigenous rights, which attracted thousands.
Controversy erupted when videos posted on social media showed some students and staff alongside pro-Palestinian protesters, and joining in chants, such as"From Turtle Island to Palestine, Occupation is a Crime."The province ordered an investigation into what happened, and the board has sincesaid students are not to attend protests or rallies during school hours -- rules many boards across the province already have in place.
The TDSB, in its Sept. 11 letter to Education Minister Jill Dunlop, said there ' s an"urgent need for a comprehensive and unified approach to addressing hate and geopolitical tensions within Ontario ' s schools."ARTICLE CONTINUES BELOW ARTICLE CONTINUES BELOW It asks the province to develop a clear definition of hate that tackles growing antisemitism and other forms of hate ; create a policy that will help boards manage student and staff activities related to broader geopolitical conflicts ; and provide resources and @ @ @ @ @ @ @ @ @ @.
In April, the TDSB asked the province to establish a new data collection tool for incidents of hate and racism, so boards in Ontario could identify \textbf{trends}, target interventions and track how issues were resolved.
The TDSB collects data on racism, bias and hate, but Shan says it can't analyze the data in a sophisticated way.
Last month, the TDSB also wrote to the Ontario College of Teachers asking it to update some professional advisories, including some guidance for teachers to better support students stressed by world conflicts, as well as clarifying for educators what is acceptable to post on social media and when sharing commentary or personal views violates professional expectations.
The college said it is working on the concerns it has heard from boards, and expects to issue an advisory on hate and discrimination early in 2025.
Ahead of Oct. 7, Dunlop sent a memo to boards telling them to"be vigilant in ensuring classrooms remain safe, inclusive and welcoming for all students and staff. @ @ @ @ @ @ @ @ @ @ political opinions, they are not entitled to disseminate political biases into our classrooms."Kristin Rushowy is a Toronto-based reporter covering Ontario politics for the Star.
Follow her on Twitter : @krushowy.
Isabel Teotonio is a Toronto-based reporter covering education for the Star.
Follow her on Twitter : @Izzy74.

% matched lemmas: coverage, exclusive, necessarily, trend
\end{textsample}
